\section{Panorama global de las empresas de modelos y de productos de IA}

Tras el análisis de las barreras de entrada de los productos, esta sección aborda el panorama de las empresas. En el mapa de empresas del Q1 de 2025, OpenAI, Anthropic, Google, DeepSeek, ByteDance, Mira/Thinking Machines, SSI, Cursor, Manus y otros actores aparecen en un mismo mapa competitivo. Las opiniones de Guangmi (广密) son marcadamente subjetivas, pero ayudan a entender la posición de los distintos tipos de empresa en la cadena de valor de la IA.

\begin{figure}[H]
\centering
\includegraphics[width=0.92\textwidth]{model-company-portfolio.png}
\caption{Panorama global de las empresas de modelos: cada empresa apuesta de forma distinta por el modelo, el producto, el ecosistema, el código abierto y la estructura de capital. Diagrama conceptual de elaboración propia, basado en la conversación de 01:15:11--01:54:32.}
\end{figure}

\begin{knowledgebox}{Lectura de la figura: no todas las empresas de IA compiten en la misma categoría}
Anthropic se orienta más a empresas y a Coding; OpenAI apuesta a la vez por el modelo y por el producto; ByteDance tiene productos y tráfico; DeepSeek representa la eficiencia del código abierto; Cursor y Manus forman la capa de ejecución de aplicaciones. Las condiciones de éxito de cada posición son completamente distintas, por lo que no basta con una sola clasificación de benchmark.
\end{knowledgebox}

\subsection{GPT-4.5, GPT-5 y los riesgos de OpenAI}

Esta sección repasa la conversación sobre OpenAI. La entrevista aborda si GPT-4.5 va a la delantera, por qué se retrasó GPT-5, si OpenAI corre riesgo de fracasar, su relación con Microsoft y su apoyo al protocolo MCP de Anthropic. Al organizar este material como curso, conviene verlo como un conjunto de variables: el ritmo del modelo base, el ritmo del reasoning model, el crecimiento de usuarios del producto, la estructura de la alianza en la nube, el ecosistema de protocolos y la estabilidad organizacional.

\begin{warningbox}{No juzgues a una empresa por un solo lanzamiento}
Tanto los riesgos como las ventajas de OpenAI son grandes: tiene la marca más fuerte, el principal punto de acceso de usuarios y el mejor posicionamiento de producto en la mente del público, pero también enfrenta una creciente complejidad organizacional, su relación con los proveedores de nube, el ritmo de la investigación y la presión de comercialización. Una sola versión de un modelo no decide quién gana a largo plazo.
\end{warningbox}

\subsection{DeepSeek, la eficiencia del código abierto y la señal de China}

Esta sección se ocupa de DeepSeek. La entrevista coloca a DeepSeek como la estrella del Q1 y subraya su impacto en el panorama global de los modelos y en el discurso sobre la eficiencia del código abierto. Su valor no se reduce a que «el modelo es barato»: demuestra que los equipos chinos pueden influir en la percepción global mediante la eficiencia de ingeniería, las decisiones de arquitectura, la recipe de entrenamiento y la difusión en código abierto. Guangmi incluso usa el lenguaje de un portafolio de inversión para expresar que le daría un peso elevado a DeepSeek, lo que refleja la importancia que concede a la vía de la eficiencia del código abierto.

\begin{importantbox}{El significado de la señal de DeepSeek}
DeepSeek llevó a la industria a reevaluar la posición de los equipos chinos en modelos frontier, eficiencia de entrenamiento, ecosistema de código abierto e influencia sobre los desarrolladores de todo el mundo. No es solo una historia de precios, sino una historia en la que se suman capacidad, costo, apertura y difusión.
\end{importantbox}

\subsection{Manus, Perplexity y Cursor: empresas de producto con gran capacidad de ejecución}

Esta sección se ocupa de las empresas de aplicaciones. En la entrevista, a Manus, Perplexity y Cursor se les llama «los que le roban el fuego a los modelos» e incluso, en tono de broma, «los reyes del envoltorio». La expresión parece mordaz, pero lo que en realidad destaca es su capacidad de ejecución: convierten con rapidez en experiencia de usuario las capacidades que las empresas de modelos aún no han convertido en producto. Perplexity reorganizó la búsqueda y la investigación, Cursor convirtió el Coding en un flujo de trabajo de uso frecuente y Manus llevó ante el público la experiencia de un Agent de propósito general.

\begin{knowledgebox}{Términos clave: empresas de modelos, empresas de aplicaciones y empresas de infraestructura}
\noindent\begin{tabularx}{\textwidth}{p{0.20\textwidth}p{0.34\textwidth}X}
\toprule
Tipo de empresa & Activos centrales & Riesgos principales \\
\midrule
Empresa de modelos & base model, infra de entrenamiento, talento de investigación, API/plataforma & Alto costo de entrenamiento, organización compleja, fuerte presión por el producto. \\
Empresa de aplicaciones & Punto de acceso de usuarios, flujos de trabajo, estado de las tareas, datos de retroalimentación & Que las empresas de modelos la copien o que los costos la ahoguen. \\
Empresa de infraestructura & Inferencia, entrenamiento, entornos, protocolos de herramientas, recursos en la nube & Debe adaptarse a los cambios en la demanda de los modelos y las aplicaciones. \\
Empresa de modelos de código abierto & Eficiencia de costos, comunidad, transparencia, difusión del ecosistema & Presión por la comercialización y por financiar el entrenamiento continuo. \\
\bottomrule
\end{tabularx}
\end{knowledgebox}

\subsection{MCP, protocolos y nichos del ecosistema}

Esta sección completa la capa de protocolos. La entrevista menciona que OpenAI apoya el protocolo MCP de Anthropic, lo que muestra que en el ecosistema de Agents no solo se compite por el modelo, sino también por la forma de conectar herramientas, la forma de transmitir el contexto y los límites de seguridad. Si un protocolo se convierte en estándar de facto, conectará los modelos, las herramientas, los sistemas empresariales y el ecosistema de desarrolladores. Las empresas de modelos deben evitar que el protocolo convierta su modelo en una mercancía intercambiable y, al mismo tiempo, evitar quedar aisladas del ecosistema.

\begin{importantbox}{La siguiente capa de competencia en el ecosistema de Agents son los protocolos}
Navegadores, IDE, SaaS empresariales, sistemas de archivos y bases de datos pueden convertirse en herramientas de un Agent. Quien defina cómo el modelo lee el contexto, invoca herramientas y devuelve resultados de forma segura estará definiendo las reglas de tránsito del ecosistema de Agents.
\end{importantbox}

\subsection{Resumen de la sección}

El panorama global de las empresas de IA no puede ordenarse solo según «quién tiene el modelo más fuerte». El modelo base, la eficiencia del código abierto, la ejecución de aplicaciones, el ecosistema de protocolos, las alianzas en la nube y la estructura de capital están reacomodando las posiciones. El cambio real del Q1 de 2025 es que la capacidad de los modelos empieza a dirigirse con más claridad hacia el Coding, los Agents y los productos de flujo de trabajo.
